\section{Angelic Knowledge}\label{sec:st-knowledge}

How does a creature with no eyes, no ears, and no imagination recognize the exact face of the child it guards, yet remain utterly blind to the silent prayer turning inside that child's heart? In Questions~54 through~58 of the \emph{Prima Pars}, St.~Thomas Aquinas maps an intelligence wholly unlike human reasoning. An angelic mind never collects impressions from the senses, never abstracts universals from phantasms, and never works step by step from premises to conclusions. Yet even at the summit of created intellect, angelic knowledge stops abruptly at three frontiers reserved to God alone: the uncreated future, the sanctuary of the free heart, and the hidden counsels of grace.

\subsection{The Cognitive Power of the Angels (\ST{I, q.~54})}

Aquinas opens his account of angelic cognition by safeguarding the boundary between created intelligence and the Creator. First, an angel's act of understanding cannot be its substance (\ST{I, q.~54, a.~1}). Action is the actuality of an operative power just as existence (\emph{esse}) is the actuality of essence; because no creature's existence is identical with its substance, still less can its action be its substance. Were an angel's act of understanding identical with its substance, there could be only one subsisting act of understanding---God Himself---and there could be no degrees of more and less perfect intelligence among angels. Aristotle and Averroes speak of the active intellect as identical with its act only concomitantly, because it is perpetually in act (\ST{I, q.~54, a.~1, ad~1}); Augustine's identification of ``to live'' with ``to be'' distinguishes concrete existence from the abstract vital operation (\ST{I, q.~54, a.~1, ad~2}); and an immanent act of knowing requires no middle barrier between knower and known, but follows upon their union (\ST{I, q.~54, a.~1, ad~3}).

Nor is understanding identical with existing in an angel (\ST{I, q.~54, a.~2}). Adopting Aristotle's distinction (\emph{Metaphysics}~IX) between transient action that passes outward into another subject (such as burning or cutting) and immanent action that abides within the agent (such as sensing, understanding, and willing), Aquinas notes that immanent intellectual operations possess a \emph{simple infinity} in their formal objects: understanding reaches universal truth and willing reaches universal good, both convertible with being itself. Every creature's existence, by contrast, is contracted to one finite genus and species. God alone, whose being is simply infinite and comprehends all things in Himself (Pseudo-Dionysius, \emph{De divinis nominibus}~V), is His own understanding and His own willing; as Dionysius remarks elsewhere (\emph{De divinis nominibus}~IV), an angel's understanding is its ``movement,'' whereas its substantial existence is not movement (\ST{I, q.~54, a.~2}).

A third real distinction follows: an angel's intellective power is not identical with its essence (\ST{I, q.~54, a.~3}). Pseudo-Dionysius writes that ``the angels are divided into substance, power, and operation'' (\emph{Celestial Hierarchy}~XI). Because every potency is proportioned to its proper act, essence is ordered to substantial existence while operative power is ordered to accidental operation; wherever existence and operation differ, essence and intellective power must differ as well. Dionysius calls angels ``intellects'' and ``minds'' not because their faculty equals their essence, but because all their cognition is purely intellectual (\ST{I, q.~54, a.~3, ad~1}). Boethius's axiom that a simple form cannot be a subject (\emph{De Trinitate}~I) holds strictly of God, who is Pure Act; a created subsisting form composed of essence and existence readily serves as the subject of proper accidents flowing from its species, contrasting prime matter's potency to substantial being with operative power's potency to operation (\ST{I, q.~54, a.~3, ad~2--3}). Here Franciscan theology takes another road: St.~Bonaventure and John Duns Scotus resist a real distinction between an intellectual creature's essence and its faculties, treating them instead as identical or formally distinct.

Because angels do not gather knowledge from material bodies, they possess neither an active (agent) nor a passive (possible) intellect in Aristotle's sense (\ST{I, q.~54, a.~4}). Human beings need a passive intellect because our minds begin in potentiality before understanding actually, and an active intellect to strip intelligible forms out of sensory phantasms (\emph{De anima}~III). Angels are naturally endowed with actual intelligible species from creation and possess no phantasms to abstract from; their reception of supernatural enlightenment from higher angels or God is not the natural potentiality of a human passive intellect (\ST{I, q.~54, a.~4, ad~1--2}). Consequently, angelic cognition is exclusively intellectual, with neither sense nor imagination (\ST{I, q.~54, a.~5}). As St.~Gregory the Great puts it, ``man senses in common with the brutes, and understands with the angels'' (\emph{Homiliae in Evangelia}~XXIX). Patristic expressions that attribute ``sense'' (St.~Augustine, \emph{City of God}~VIII), ``experience'' (St.~Isidore, \emph{De summo bono}), or ``perverted phantasy'' (Pseudo-Dionysius, \emph{De divinis nominibus}~IV) to angels or demons either speak provisionally within the Platonist aerial-body hypothesis or use sensory vocabulary analogically---``sensing'' for certainty, ``experience'' for grasping singulars without organs, ``memory'' in Augustine's intellectual sense (\emph{De Trinitate}~X), and ``perverted phantasy'' for a disordered practical estimate of the good (\ST{I, q.~54, a.~5, ad~1--3}).

\angeldossier{\ST{I, q.~54, aa.~1--5} (Cognitive power of the angels: distinction of understanding, \emph{esse}, and intellect from essence; absence of agent/possible intellect and sensory powers).}{Common teaching (aa.~1--2, 5); Thomist position (aa.~3--4).}{Pseudo-Dionysius, \emph{De divinis nominibus}~IV--V and \emph{De caelesti hierarchia}~XI; St.~Augustine, \emph{De Trinitate}~X, \emph{Confessions}~XII, and \emph{De civitate Dei}~VIII, XXI; St.~Gregory the Great, \emph{Homiliae in Evangelia}~XXIX; St.~Isidore of Seville, \emph{De summo bono}; Boethius, \emph{De Trinitate}~I; Aristotle, \emph{De anima}~II--III and \emph{Metaphysics}~I, IX, XII; Averroes on \emph{De anima}~III and \emph{Metaphysics}~XII.}{Franciscan theologians maintain an identity or formal distinction between an intellectual creature's essence and its cognitive faculties (St.~Bonaventure; John Duns Scotus) (a.~3); none recorded on aa.~1--2, 4--5.}

\subsection{The Medium of Angelic Knowledge (\ST{I, q.~55})}

If an angel has no senses to gather concepts from the world, how do objects enter its mind? First, angels do not know all things through their own substance (\ST{I, q.~55, a.~1}). Only God's infinite essence contains the proper likeness of every creature; an angel's essence is restricted to one single species, so its intellect---open to universal truth---must be perfected by superadded intelligible species. Pseudo-Dionysius affirms that ``the angels are enlightened by the forms [\emph{rationibus}] of things'' (\emph{De divinis nominibus}~IV). When Aristotle writes that in immaterial things intellect and understood are the same (\emph{Metaphysics}~XII; \emph{De anima}~III), he means the intellect in act informed by its intelligible species, not an identity between the angel's finite essence and every object it knows; other creatures appear in an angel's own essence only under a generic formality, never in their proper natures (\ST{I, q.~55, a.~1, ad~1--3}).

Second, angels do not abstract those intelligible species from sensible things, but receive them as connatural, concreated gifts directly from God together with their intellectual nature (\ST{I, q.~55, a.~2}). Pseudo-Dionysius observes that ``the angels do not gather their Divine knowledge from things divisible or sensible'' (\emph{De divinis nominibus}~VII). Aquinas illustrates this by comparing the hierarchies of bodies and spirits: heavenly bodies are naturally complete in form, whereas earthly matter acquires forms successively; likewise, human souls---naturally united to bodies---acquire their intelligible perfection bit by bit through bodily senses, while pure spirits receive their full natural complement of species at creation by an intelligible outpouring from God. St.~Augustine expresses the same order in \emph{De Genesi ad litteram}~(II.8): ``The other things which are lower than the angels are so created that they first receive existence in the knowledge of the rational creature, and then in their own nature.'' Lacking imagination, angels could not abstract from matter even if they tried, nor do they travel locally to learn what happens on earth; they move locally to operate, not to acquire information (\ST{I, q.~55, a.~2, ad~2--3}). John Duns Scotus later departed from Aquinas on this point, holding that angels can also acquire new intelligible species directly from created objects alongside their innate species.

Third, the higher an angel stands in nature, the fewer and more universal are the intelligible species by which it knows the entire cosmos (\ST{I, q.~55, a.~3}). God knows everything through one utterly simple essence; created intelligences approach that simplicity by degrees, grasping in fewer, richer forms what lower minds must split across many narrower concepts. Pseudo-Dionysius testifies that higher angels share a more universal knowledge (\emph{Celestial Hierarchy}~XII), and the \emph{Liber de causis} teaches that higher intelligences hold more universal forms. We glimpse a faint shadow of this even among men, where a slow learner needs every detail spelled out separately while a sharp mind grasps a wide range of truths from a few principles (\ST{I, q.~55, a.~3}). In his earlier \emph{De veritate} (q.~8, a.~10), Aquinas offers an architectonic analogy: a master builder (\emph{architectus}) directs an entire building through one unified form of his art, whereas subordinate stonemasons and woodcutters must each be guided by separate, specialized crafts.

A crucial distinction prevents misunderstanding (\ST{I, q.~55, a.~3, ad~2--3}). Knowing ``universally'' can refer to the \emph{object} known---knowing a man merely as ``animal''---which is vague and imperfect. Angelic universality lies instead on the side of the \emph{medium}: one eminent, richly concentrated form contains distinctly every proper detail of the realities beneath it, much as human universal prudence encompasses in one habit the particular instincts of the lion's magnanimity and the fox's wariness, and as the Divine essence eminently serves as the proper likeness of every individual creature.

\angeldossier{\ST{I, q.~55, aa.~1--3} (Medium of angelic knowledge: intelligible species distinct from substance, innate reception from God, and fewer, more universal species in higher angels).}{Common teaching (a.~1); Thomist position (aa.~2--3).}{Pseudo-Dionysius, \emph{De divinis nominibus}~IV, VII and \emph{De caelesti hierarchia}~XII; St.~Augustine, \emph{De Genesi ad litteram}~II.8; \emph{Liber de causis}; Aristotle, \emph{De anima}~II--III and \emph{Metaphysics}~XII; St.~Thomas Aquinas, \emph{De veritate}~q.~8, aa.~8--10.}{John Duns Scotus admits intelligible species acquired from created objects alongside innate species (a.~2); none recorded on aa.~1, 3.}

\subsection{Angelic Knowledge of Immaterial Things (\ST{I, q.~56})}

Turning from the medium of knowledge to its objects, Aquinas considers first how an angel knows immaterial beings: itself, other angels, and God. An angel knows itself directly through its own substance, without any intermediary species (\ST{I, q.~56, a.~1}). Because an angel is an immaterial subsisting form, it is actually intelligible in itself and immediately united to its own intellect; as St.~Augustine writes, ``the angel knew himself when he was established, that is, enlightened by truth'' (\emph{De Genesi ad litteram}~II). Singularity never obstructs intelligibility---only matter does---so an immaterial singular is immediately transparent to itself (\ST{I, q.~56, a.~1, ad~2}). Noting a textual divergence between two Latin versions of Pseudo-Dionysius (\emph{Celestial Hierarchy}~VI), where the older version read ``they do not know their own powers'' and the newer read ``they knew their own powers,'' Aquinas adds that even the older rendering remains true if read as denying comprehensive knowledge of how angelic power flows from Divine Wisdom (\ST{I, q.~56, a.~1, ad~1}).

How does one angel know another (\ST{I, q.~56, a.~2})? Not through its own essence, which is specifically distinct from every other angel; and not directly through the other angel's essence, a thesis proposed by Averroes and several scholastics, because only the Triune God indwells a creature's essence directly (\emph{De veritate}~q.~8, a.~7, where Aquinas also rejects a literalist reading of Dionysius's figure of higher angels as ``food'' for lower). Instead, one angel knows another through concreated intelligible species impressed by the Word of God at creation. From the Word, the archetypes of all creatures proceed both into natural subsistence and into angelic intellects: each angel receives its own form naturally and intelligibly, and every other angel's form intelligibly. Aquinas distinguishes sharply between \emph{esse naturale} (``natural being'') and \emph{esse intentionale} (``intentional'' or intelligible being): an angel is a subsisting form in natural being, while its likeness in another angel's mind possesses intentional being, like color passing through illuminated air compared with pigment on a wall (\ST{I, q.~56, a.~2, ad~3}). Were God to create a new angel today, He would simultaneously impress its species in existing angels, just as a builder widens foundations when adding to a house (\ST{I, q.~56, a.~2, ad~4}).

By his own natural principles, an angel also knows God (\ST{I, q.~56, a.~3}). If human beings attain natural knowledge of God from created things (Rom~1:19--20), angels, far mightier in intelligence, know Him even more clearly. Yet no created intellect---human or angelic---can behold God's essence directly by natural power alone (\ST{I, q.~12, a.~4}; \emph{De veritate}~q.~8, aa.~1--3). Aquinas lays out three modes of knowing an object: first, by the presence of its essence in the knower, as light is seen in the eye (beatific vision of God, or angelic self-knowledge); second, by a likeness drawn directly from the object itself, as the eye sees a stone; third, in a mirror, as earthly men know God through creatures (1~Cor~13:12). Here Aquinas subtly refines his earlier teaching. In \emph{De veritate} (q.~8, a.~3, ad~17), he placed angelic natural knowledge squarely in the \emph{second} mode, via an impressed likeness strictly midway between essential vision and human specular knowledge: \emph{secundo modo videre Deum est naturale Angelo; tertio autem modo videre Deum est naturale ipsi homini}, ``to see God in the second way is natural to the angel; to see God in the third way is natural to man.'' In \ST{I, q.~56, a.~3}, Aquinas brings the image even closer to St.~Paul's mirror: because God's image is impressed upon the angel's very essence, yet no created likeness suffices to represent the divine essence as it is, ``such knowledge then approaches rather to the specular kind; because the angelic nature is itself a kind of mirror representing the Divine image.'' Pseudo-Dionysius's apophatic declarations that God transcends all knowledge (\emph{De divinis nominibus}~I, VII) exclude natural vision of God's essence and comprehension, not genuine natural knowledge (\ST{I, q.~56, a.~3, ad~1--2}).

\angeldossier{\ST{I, q.~56, aa.~1--3} (Angelic knowledge of immaterial beings: direct self-knowledge, knowledge of other angels by innate species, and natural specular knowledge of God).}{Common teaching (aa.~1, 3); Thomist position (a.~2).}{Rom~1:19--20; 1~Cor~13:12; Pseudo-Dionysius, \emph{De caelesti hierarchia}~VI and \emph{De divinis nominibus}~I, VII; St.~Augustine, \emph{De Genesi ad litteram}~II; \emph{Liber de causis}; Aristotle, \emph{De anima}~III; St.~Thomas Aquinas, \emph{De veritate}~q.~8, aa.~3, 6--7.}{Averroes and several scholastics proposed direct intuition of another angel's essence without mediating species (a.~2); in the earlier \emph{De veritate}~q.~8, a.~3, ad~17, Aquinas classifies natural angelic knowledge of God as strictly intermediate by impressed species rather than specular (a.~3); none recorded on a.~1.}

\subsection{Material Things, Singulars, and the Three Boundaries of Knowledge (\ST{I, q.~57})}

An angel knows not only spirits, but the entire material creation (\ST{I, q.~57, a.~1}). Whatever a lower power grasps, a higher power grasps in a nobler mode; by the Dionysian principle of hierarchy (\emph{De divinis nominibus}~I; \emph{Celestial Hierarchy}~IV), realities that subsist imperfectly in lower beings pre-exist more simply and immaterially in higher ones---supersubstantially in God and intelligibly in angels. Material bodies do not lower the angelic intellect, because the perfection of the knower is the immaterial species inside its mind, not the heavy matter outside it (\ST{I, q.~57, a.~1, ad~1--2}).

Far more searching is the question of \emph{singulars}: how does a pure spirit without eyes know this particular stone, this child, this street (\ST{I, q.~57, a.~2})? Aquinas anchors his argument directly in angelic guardianship: no one can guard what he does not know, yet Scripture promises, ``For he hath given his angels charge over thee; to keep thee in all thy ways'' (Ps~90:11 [Heb~91:11]). He rejects two philosophical errors. To deny angels all knowledge of singulars contradicts Catholic faith in angelic ministry (Heb~1:14; Eccles~5:5) and even Aristotelian cosmology. Avicenna's compromise---allowing separate intelligences to know singulars only through their universal causes, the way an astronomer calculates a coming eclipse from orbital tables---likewise fails: the astronomer knows the eclipse universally, not as happening \emph{here and now}, unless he looks up with his eyes, whereas providence and guardianship govern individuals here and now. Just as human common sense knows in one faculty everything five separate outer senses perceive, an angel's single superior intellect knows both universals and singulars; Aristotle himself judged it absurd that something known to men should remain unknown to higher intellect (\emph{De anima}~I).

How, then, can an immaterial universal species represent a material individual? Aquinas solves the puzzle from the direction of causality (\ST{I, q.~57, a.~2}; \emph{De veritate}~q.~8, a.~11). Human concepts are abstracted \emph{from} sensible things, shedding individuating matter along the way, so our concepts grasp only universal natures. Angelic species flow in the opposite direction---downward from the creative essence of God. Because God causes the entire substance of every creature, both its universal form and its individuating matter, God's creative knowledge reaches every singular in its concrete individuality; and the innate species imparted to angels, as manifold likenesses of that creative Divine essence, represent creatures both in their universal species and in their individual conditions (\ST{I, q.~57, a.~2, ad~2--3}).

Yet this luminous gaze meets three inviolable walls. The first is the \emph{future} (\ST{I, q.~57, a.~3}). Isaias makes foreknowledge of what is to come the exclusive proof of divinity: ``Shew the things that are to come hereafter, and we shall know that ye are gods'' (Is~41:23). Angels know future events \emph{in their causes}: with certainty when an effect follows necessarily (such as tomorrow's sunrise), and conjecturally---far more reliably than the sharpest physician predicting a patient's recovery---when causes produce effects in most cases. Contingent, chance, or free future events they cannot know in themselves, because future things do not yet possess actual existence conformed to their innate species (\ST{I, q.~57, a.~3, ad~3}). Only God beholds all future events as present in His simple eternity, whereas angelic cognition unfolds within a discrete succession of concepts---``God moves the spiritual creature according to time'' (St.~Augustine, \emph{De Genesi ad litteram}~VIII; \ST{I, q.~57, a.~3, ad~2}). In \emph{De malo} (q.~16, a.~7), Aquinas crystallizes this contrast with the image of a watchman stationed on a high watchtower (\emph{in alta specula}): from that height he sees at once every traveler passing along the road below, whereas anyone walking on the road itself inside the line of travelers sees only those beside or behind him. Except when God reveals a future event or an angel acts upon its own intention, the uncreated future stays sealed to every created spirit.

The second wall is the \emph{secret of the human heart} (\ST{I, q.~57, a.~4}). Jeremias proclaims: ``The heart is perverse above all things, and unsearchable, who can know it? I am the Lord who search the heart'' (Jer~17:9--10), and St.~Paul asks, ``For what man knoweth the things of a man, but the spirit of a man that is in him?'' (1~Cor~2:11). Indirectly, both good angels and demons read subtle bodily clues---flushing, pulse, posture, humors, imaginative stirrings---with vastly greater acuity than any physician (St.~Augustine, \emph{De divinatione daemonum}, qualified cautiously in \emph{Retractationes}~II.30). Directly in themselves, however, thoughts of the mind and affections of the will lie open to God alone. In \emph{De malo} (q.~16, a.~8), Aquinas sharpens the psychological proof: a superior spiritual intellect can even perceive the habitual \emph{intelligible species} stored inside a human mind, yet the actual \emph{use of the species} (\emph{usus speciei}, actual thinking) hangs solely on the free will. Because a rational creature's will is immediately subordinate only to God as its Supreme Good and last end, no angel or demon can read what a person freely chooses to think or love unless that person wills to disclose it---just as a king's direct order to a citizen stays hidden from a provincial bailiff. Even after the resurrection, bodily splendor will manifest the saints' degree of glory without stripping away the will's interior privacy (\ST{I, q.~57, a.~4, ad~1--3}).

The third wall is the \emph{mysteries of grace} (\ST{I, q.~57, a.~5}). In Isaias~63:1, even the loftiest spirits ask in wonder, ``Who is this that cometh from Edom?''---a text Pseudo-Dionysius cites to show that the highest angels learn divine counsels from Jesus Himself (\emph{Celestial Hierarchy}~VII). Because mysteries of grace flow purely from God's free will, no angel can reach them by natural intellect (1~Cor~2:11). In the beatific vision of the Word, blessed angels know mysteries of grace progressively and unequally according to God's revelation (1~Cor~2:10), higher angels seeing deeper counsels and illuminating lower ones; for beholding God's essence does not mean comprehending His infinite wisdom (\emph{De veritate}~q.~8, aa.~2, 4; \ST{I, q.~57, a.~5, ad~2}). All blessed angels received knowledge of the Incarnation in general from the beginning of their beatitude, since all angelic ministry serves human salvation (Heb~1:14); its concrete timing and circumstances they learned as salvation history unfolded (\ST{I, q.~57, a.~5, ad~1, 3}). Commenting on Ephesians~3:10---``that the manifold wisdom of God may be made known to the principalities and powers in heavenly places through the church''---Aquinas resolves how blessed angels learn through earthly apostles without reversing the heavenly hierarchy: \emph{non quod angeli hoc didicerint ab eis, sed in eis}, ``not that the angels learned this from them, but in them,'' recognizing God's hidden design as it comes to pass in the Church, like onlookers perceiving a master builder's concealed plan once the house stands finished (\emph{Super Ephesios}~3:10).

\angeldossier{\ST{I, q.~57, aa.~1--5} (Angelic knowledge of material things, singulars, future contingents, secret thoughts of hearts, and mysteries of grace).}{Common teaching (aa.~1--5).}{Ps~90:11 [Heb~91:11]; Ps~118:100 [Heb~119:100]; Job~28:17; Eccles~5:5; Is~41:23; Is~63:1; Jer~17:9--10; Amos~3:7; Matt~22:30; 1~Cor~2:10--11; 2~Cor~12:2; Eph~3:4--5, 10; 1~Tim~3:16; Heb~1:14; Pseudo-Dionysius, \emph{De divinis nominibus}~I and \emph{De caelesti hierarchia}~IV, VII; St.~Augustine, \emph{De Genesi ad litteram}~V.19, VIII, XII.28, \emph{De divinatione daemonum}, and \emph{Retractationes}~II.30; St.~Gregory the Great, \emph{Moralia in Iob}~XVIII and \emph{Homiliae in Ezechielem}~XVI; \emph{Liber de causis}; Aristotle, \emph{De anima}~I, III, \emph{Posterior Analytics}~I, \emph{Metaphysics}, and \emph{Nicomachean Ethics}~III.12; St.~Thomas Aquinas, \emph{De veritate}~q.~8, aa.~4, 11--13, \emph{De malo}~q.~16, aa.~7--8, and \emph{Super Ephesios}~3:10.}{Avicenna and some Aristotelians restricted separate intelligences to universal knowledge, denying direct cognition of material singulars (a.~2); none recorded on aa.~1, 3--5.}

\subsection{The Mode of Angelic Knowledge (\ST{I, q.~58})}

Having established what angels know, Aquinas examines \emph{how} their intellect operates. First, is an angelic intellect ever in potentiality (\ST{I, q.~58, a.~1})? Applying Aristotle's two levels of potentiality (\emph{De anima}~III), Aquinas distinguishes sharply. In first potentiality---lacking habitual species---an angel is \emph{never} in potentiality regarding natural knowledge, though he remains open to fresh supernatural revelations. In second potentiality---possessing habitual species without currently thinking of every one of them---an angel can be in potentiality regarding natural knowledge, yet \emph{never} regarding the beatific vision of the Word, which is ceaseless act (St.~Augustine, \emph{De Genesi ad litteram}~II). When St.~Peter speaks of realities ``on whom the angels desire to look'' (1~Pet~1:12), that desire excludes weariness in vision or points to new revelations ordered to their ministry (\ST{I, q.~58, a.~1, ad~2}).

Second, can an angel understand many things at the exact same time (\ST{I, q.~58, a.~2})? St.~Augustine praises the scope of angelic comprehension (\emph{De Genesi ad litteram}~IV.32) and compares beatific simultaneity to surveying all knowledge in one single glance rather than flitting from thought to thought (\emph{De Trinitate}~XV.16). Aquinas introduces a clean metaphysical rule: unity of operation demands unity of informing species, just as one piece of wax cannot hold two shapes simultaneously (\emph{De veritate}~q.~8, a.~14). An angel therefore understands simultaneously \emph{everything} grasped through one intelligible species---including all things it beholds in the single Divine essence, and all creatures unified within one comprehensive innate species---but does not actuate several distinct innate species at the same instant (\ST{I, q.~58, a.~2}). Later Franciscan and Jesuit masters, notably John Duns Scotus and Francisco Su\'arez, rejected this restriction and permitted simultaneous actuation of multiple distinct species.

Third and fourth, angelic knowledge is neither \emph{discursive} (\ST{I, q.~58, a.~3}) nor \emph{compositional} (\ST{I, q.~58, a.~4}). Aquinas draws the defining line between \emph{intellectus} (``intellect'') and \emph{ratio} (``reason''). Human beings are called \emph{rational} creatures because our feeble intellectual light forces us to reason discursively, advancing step by step from known principles to previously unseen conclusions, and composing or dividing predicates against subjects to build definitions. Only in our immediate, effortless grasp of first principles do we touch, as Pseudo-Dionysius says, the lowest hem of angelic life (\emph{De veritate}~q.~8, a.~15; \emph{De divinis nominibus}~VII). Angels are properly called \emph{intellectual} beings because their full light beholds at a single glance every conclusion virtually contained in a principle and every attribute belonging to or excluded from a thing's essence:
\begin{quotation}
Hence, since the intellectual light is perfect in the angel, for he is a pure and most clear mirror, as Dionysius says (\emph{Div.~Nom.}~iv), it follows that as the angel does not understand by reasoning, so neither does he by composing and dividing. Nevertheless, he understands the composition and the division of enunciations, just as he apprehends the reasoning of syllogisms: for he understands simply, such things as are composite, things movable immovably, and material things immaterially. (St.~Thomas Aquinas, \ST{I, q.~58, a.~4})
\end{quotation}
Thus angels understand human arguments and converse with men in affirmative and negative sentences without ever needing syllogisms to discover truth (\ST{I, q.~58, a.~3, ad~1--3}; \ST{a.~4, ad~3}).

This intuitive simplicity immediately illuminates the question of falsehood (\ST{I, q.~58, a.~5}). Aristotle (\emph{De anima}~III; \emph{Metaphysics}~IX) and St.~Augustine (\emph{De diversis quaestionibus LXXXIII}, q.~32) agree that simple apprehension of what a thing is never errs; human error enters when we misjoin terms by composition and division, as in imagining a ``four-footed flying beast.'' Consequently, no angel can ever err \emph{per se}. In good angels, not even accidental error can occur: regarding unrevealed supernatural decrees they have only simple \emph{nescience} (Pseudo-Dionysius, \emph{Ecclesiastical Hierarchy}~VI), and because their upright wills subject every judgment to divine wisdom, they never assess a creature outside God's ordinance. Demons, too, remain infallible regarding natural properties; yet because their rebellious wills reject divine wisdom, they fall into accidental error in supernatural matters---seeing a dead man and assuming he will not rise again, or beholding Christ and judging Him not to be God (\ST{I, q.~58, a.~5}; cf.~\emph{De malo}~q.~16, a.~6).

Question~58 closes with St.~Augustine's celebrated distinction between \emph{morning} and \emph{evening} knowledge (\ST{I, q.~58, aa.~6--7}). Reading the six days of Genesis~1 not as solar circuits (since the sun appeared only on the fourth day) but as the single day of angelic knowledge directed to six classes of creatures (\emph{De Genesi ad litteram}~IV.22--24, 31), Augustine contrasts the two gazes in \emph{City of God}~(XI.29):
\begin{quotation}
They know the creature also, not in itself, but by this better way, in the wisdom of God, as if in the art by which it was created; and, consequently, they know themselves better in God than in themselves\dots{} In Him, therefore, they have, as it were, a noonday knowledge; in themselves, a twilight knowledge. (St.~Augustine, \emph{City of God}~XI.29)
\end{quotation}
\emph{Morning knowledge} (\emph{cognitio matutina}) beholds creatures in their primordial being in the Word, the source from which all being flows; \emph{evening knowledge} (\emph{cognitio vespertina}) beholds creatures in the created being they possess in their own proper natures (\ST{I, q.~58, a.~6}). For the holy angels, evening never darkens into night, because they immediately refer their knowledge of creatures back to the praise of God, turning evening into a new morning (\emph{De Genesi ad litteram}~IV.31.48); only the demons cling to creatures apart from God and plunge into night (\ST{I, q.~58, a.~6, ad~2}).

Are morning and evening knowledge one cognitive act or two (\ST{I, q.~58, a.~7})? First, Aquinas rules out any misreading of the preposition ``in'': knowing things ``in their proper nature'' names the \emph{object} known, never an abstraction drawn \emph{from} creatures. Second, Aquinas refines his own earlier Parisian teaching. In \emph{De veritate} (q.~8, aa.~16--17), he explicitly refused to separate morning and evening knowledge merely by object and bound them exclusively to uncreated versus created \emph{medium} (vision in the Word versus innate species). In \ST{I, q.~58, a.~7}, he accommodates both senses gracefully: insofar as angels see the proper natural being of creatures by gazing into the Word, morning and evening knowledge are essentially one act differing only in object; insofar as they know creatures' proper nature through their innate species, morning and evening knowledge are two distinct acts, coexisting harmoniously in the blessed because evening knowledge is ordered to morning knowledge---and glory abolishes only opposing imperfections such as faith (1~Cor~13:10), not a compatible subordinate medium (\ST{I, q.~58, a.~7, ad~2--3}).

\angeldossier{\ST{I, q.~58, aa.~1--7} (Mode of angelic knowledge: act and potentiality, simultaneity under one species, non-discursion, simple apprehension without composition, immunity from error, and morning and evening knowledge).}{Common teaching (aa.~1, 3, 5--6); Thomist position (aa.~2, 4, 7).}{Gen~1:5; Matt~18:10; 1~Cor~13:10; Eph~5:8; 1~Pet~1:12; 2~Pet~1:19; Pseudo-Dionysius, \emph{De divinis nominibus}~IV, VII and \emph{De ecclesiastica hierarchia}~VI; St.~Augustine, \emph{De Genesi ad litteram}~II, IV.22--24, 31--32, \emph{De civitate Dei}~XI.29, \emph{De Trinitate}~XV.16, and \emph{De diversis quaestionibus LXXXIII}~q.~32; St.~Isidore of Seville, \emph{De summo bono}~I.10; \emph{Liber de causis}; Aristotle, \emph{Physics}~III, VIII, \emph{Topics}~II.4, \emph{Posterior Analytics}~II, \emph{De anima}~III, \emph{Metaphysics}~VII, IX, and \emph{Nicomachean Ethics}~I.8; St.~Thomas Aquinas, \emph{De veritate}~q.~8, aa.~14--17.}{John Duns Scotus and Francisco Su\'arez permit simultaneous actuation of several distinct intelligible species in an angelic intellect (a.~2); in \emph{De veritate}~q.~8, a.~16, Aquinas separates morning and evening knowledge strictly by uncreated versus created medium rather than accommodating both distinctions (a.~7); none recorded on aa.~1, 3--6.}
